%% kapa: exam=1 round=29 year=2018 period=1 form=A subject=real_estate subject_ko=부동산학원론 seq=8 number=88
%% source: Q-Net 감정평가사 기출문제 (cst00309) artlSeq=5202808 file=extracted/1008/5202808/2018년 제29회 감정평가사 1차시험 기출문제/2018년 제29회 감정평가사 1차시험 기출문제/1교시 A형.pdf
%% conversion: kordoc 4.24.51 -> tools/gen1.py (auto); review=proofread-sonnet; flags=none
\begin{question}{88}
전국에 세 개의 지역(A, B, C)과 세 개의 산업(제조업, 금융업, 숙박업)만 존재한다고 가정할 때 입지계수에 관한 설명으로 옳은 것은?
\par
{\small\centering
\begin{tabular}{|c|c|c|c|c|}\hline
\parbox{6em}{\hfill 지 역\par\noindent 산 업} & A & B & C & 전국\\\hline
\makecell{제조업\\고용자수(명)} & 150 & 170 & 195 & 515\\\hline
\makecell{금융업\\고용자수(명)} & 200 & 180 & 190 & 570\\\hline
\makecell{숙박업\\고용자수(명)} & 180 & 190 & 200 & 570\\\hline
합계(명) & 530 & 540 & 585 & 1,655\\\hline
\end{tabular}\par}
\begin{choices}
  \item B지역의 제조업은 A지역의 숙박업보다 입지계수가 낮다.
  \item A지역의 숙박업은 C지역의 금융업보다 입지계수가 높다.
  \item A지역의 숙박업과 B지역의 제조업의 입지계수는 같다.
  \item A지역의 제조업은 C지역의 숙박업보다 입지계수가 높다.
  \item B지역의 제조업은 C지역의 금융업보다 입지계수가 낮다.
\end{choices}
\end{question}
\answer{2}
